\setcounter{section}{3}

\Needspace{5\baselineskip}
\hypertarget{ux53c2ux6570ux521dux59cbux5316}{%
\section{参数初始化}\label{ux53c2ux6570ux521dux59cbux5316}}

\Needspace{5\baselineskip}
\hypertarget{ux4e00xavierglorotux521dux59cbux5316}{%
\subsection{一、Xavier（Glorot）初始化}\label{ux4e00xavierglorotux521dux59cbux5316}}

为了避免梯度消失或爆炸，我们希望每一层输出的方差与输入的方差大致相等。除了在每次计算结束后用批量规范化``拉回来''之外，我们还希望计算过程本身就能保证这一点。在线性网络中，如y\_j=sigma(x\_i*w\_ij)+b\_j，输入与输出方差的关系取决于参数的取值，因此参数如何初始化至关重要。

假设权重和输入独立同分布，且它们的均值都为 0。权重的方差为sigma\^{}2，这是我们要解的。输入的方差为Var(x\_i)=Var(X)。

考虑一个简单的线性层，输出y\_j=sigma(x\_i\emph{w\_ij)（忽略偏置），对任意j，Var(y\_j)=sigma(Var(x\_i}w\_ij))=n\_in\emph{Var(x\_i}w\_ij)（n\_in为输入数量）

\Needspace{5\baselineskip}
根据方差公式：

\[
\mathrm{Var}(Z)=E[Z^2]-(E[Z])^2
\]

\Needspace{5\baselineskip}
有：

\[
E[W_iX_i]=E[W_i]E[X_i]=0\cdot 0=0
\]

\[
\mathrm{Var}(W_iX_i)
=E[(W_iX_i)^2]-0
=E[W_i^2]E[X_i^2]
\]

\Needspace{5\baselineskip}
因为 \(E[W]=0\)，所以：

\[
\mathrm{Var}(W)=E[W^2]-(E[W])^2=E[W^2]
\]

\Needspace{5\baselineskip}
同理：

\[
\mathrm{Var}(X)=E[X^2]
\]

\Needspace{5\baselineskip}
将这些代入，可以得到：

\[
\mathrm{Var}(Y)=n_{\mathrm{in}}\cdot E[W_i^2]E[X_i^2]
\]

\Needspace{5\baselineskip}
即：

\[
\mathrm{Var}(Y)=n_{\mathrm{in}}\cdot \mathrm{Var}(W)\cdot \mathrm{Var}(X)
\]

前向传播的目标是保持方差不变，即~Var(Y) = Var(X)，故Var(W)=1/n\_in

\Needspace{5\baselineskip}
同理，对于反向传播，为了保持梯度不变性，我们希望：

\[
\mathrm{Var}\left(\frac{\partial L}{\partial z^{l-1}}\right)
=\mathrm{Var}\left(\frac{\partial L}{\partial z^l}\right)
\]

\Needspace{5\baselineskip}
这里z\_l是第l层的输出（尚未经过激活函数）。接下来进行推导：

为了推导出 \(\dfrac{1}{n_{\mathrm{out}}}\)，我们使用 Xavier（Glorot）初始化的假设：激活函数是 \(\tanh\) 或 sigmoid，并且输入 \(z\) 大部分位于其线性区域（\(z\approx 0\)）。在这个区域内，\(f(z)\approx z\)，更重要的是导数 \(f'(z)\approx 1\)。同时假设权重 \(W\)、输入 \(a\) 以及梯度 \(\dfrac{\partial L}{\partial z}\) 的均值都为 0：

\[
E[W]=0,\qquad E\left[\frac{\partial L}{\partial z^l}\right]=0
\]

并假设 \(W\)、\(\dfrac{\partial L}{\partial z^l}\) 和 \(a^{l-1}\) 相互独立。

\Needspace{5\baselineskip}
分析梯度如何从第 \(l\) 层反向传播到第 \(l-1\) 层。令：

\[
z^l=W^la^{l-1}+b^l,\qquad a^l=f(z^l)
\]

\Needspace{5\baselineskip}
其中 \(W^l\) 的形状为 \((n_l,n_{l-1})\)，\(n_{l-1}\) 是第 \(l-1\) 层的神经元数量（fan-in），\(n_l\) 是第 \(l\) 层的神经元数量（fan-out）。根据链式法则：

\[
\frac{\partial L}{\partial z^{l-1}}
=\frac{\partial L}{\partial z^l}\cdot
\frac{\partial z^l}{\partial a^{l-1}}\cdot
\frac{\partial a^{l-1}}{\partial z^{l-1}}
\]

\Needspace{5\baselineskip}
写成逐元素形式：

\[
\frac{\partial L}{\partial z^{l-1}}
=\left((W^l)^\top\cdot\frac{\partial L}{\partial z^l}\right)
\odot f'(z^{l-1})
\]

\Needspace{5\baselineskip}
看 \(\dfrac{\partial L}{\partial z_j^{l-1}}\) 中任意一个神经元 \(j\)：

\[
\frac{\partial L}{\partial z_j^{l-1}}
=\left(\sum_{k=1}^{n_l}W^l_{k,j}\frac{\partial L}{\partial z_k^l}\right)\cdot f'(z_j^{l-1})
\]

这个求和是对第 \(l\) 层所有 \(k\) 个神经元进行的，因此 \(n_l\) 正是第 \(l-1\) 层的 fan-out，记作 \(n_{\mathrm{out}}\)。

反向传播中，对于前一层（第l-1层）的一个神经元j，L对它的梯度是第l层的全部（第1至n\_out个）神经元传播而来的梯度的总和。故需要Var(W)=Var(w\_k,j)=1/n\_out.

实际操作中，取二者的调和平均数2/(n\_in+n\_out)作为方差。

\Needspace{5\baselineskip}
\hypertarget{ux4e8ckaimingux521dux59cbux5316}{%
\subsection{二、Kaiming初始化}\label{ux4e8ckaimingux521dux59cbux5316}}

\Needspace{5\baselineskip}
若激活函数使用ReLU，f'(z)在z\textgreater0时为 1，在z\textless0时为 0，其方差：

\Needspace{5\baselineskip}
先使用独立随机变量乘积的方差公式：

\[
\mathrm{Var}(AB)
=\mathrm{Var}(A)\mathrm{Var}(B)
+\mathrm{Var}(A)E[B]^2
+\mathrm{Var}(B)E[A]^2
\]

\Needspace{5\baselineskip}
令 \(A=W^l\) 且 \(B=a^{l-1}\)：

\[
\begin{aligned}
\mathrm{Var}(W^la^{l-1})
&=\mathrm{Var}(W^l)\mathrm{Var}(a^{l-1})
+\mathrm{Var}(W^l)(E[a^{l-1}])^2\\
&\quad+\mathrm{Var}(a^{l-1})(E[W^l])^2
\end{aligned}
\]

\Needspace{5\baselineskip}
应用假设 \(E[W^l]=0\)：

\[
\mathrm{Var}(W^la^{l-1})
=\mathrm{Var}(W^l)\left(\mathrm{Var}(a^{l-1})+(E[a^{l-1}])^2\right)
\]

\Needspace{5\baselineskip}
根据方差定义 \(E[X^2]=\mathrm{Var}(X)+(E[X])^2\)，可将上式简化为：

\[
\mathrm{Var}(W^la^{l-1})
=\mathrm{Var}(W^l)\cdot E[(a^{l-1})^2]
\]

\Needspace{5\baselineskip}
代回 \(\mathrm{Var}(Z^l)\) 的公式，得到前向传播主公式：

\[
\mathrm{Var}(Z^l)=n_{\mathrm{in}}\cdot \mathrm{Var}(W^l)\cdot E[(a^{l-1})^2]
\]

\Needspace{5\baselineskip}
接下来需要 \(E[(a^l)^2]\) 和 \(\mathrm{Var}(Z^l)\) 之间的关系。由于 ReLU 满足 \(a^l=\max(0,Z^l)\)，有：

\[
E[(a^l)^2]=E[(\max(0,Z^l))^2]
\]

\Needspace{5\baselineskip}
假设 \(Z^l\) 是均值为 0 的对称分布，例如高斯分布，则：

\[
\begin{aligned}
E[(a^l)^2]
&=\int_{-\infty}^{\infty}(\max(0,z))^2p(z)\,dz\\
&=\int_{-\infty}^{0}0^2p(z)\,dz+\int_{0}^{\infty}z^2p(z)\,dz\\
&=\int_{0}^{\infty}z^2p(z)\,dz
\end{aligned}
\]

\Needspace{5\baselineskip}
因为 \(p(z)\) 是对称的，所以 \(\int_0^\infty z^2p(z)\,dz\) 恰好是 \(\int_{-\infty}^{\infty}z^2p(z)\,dz\) 的一半。而：

\[
E[(Z^l)^2]=\int_{-\infty}^{\infty}z^2p(z)\,dz
\]

\Needspace{5\baselineskip}
又因为 \(E[Z^l]=0\)，所以 \(\mathrm{Var}(Z^l)=E[(Z^l)^2]\)。因此：

\[
E[(a^l)^2]=\frac{1}{2}\mathrm{Var}(Z^l)
\]

\Needspace{5\baselineskip}
将这个关系代入前向传播主公式，并要求信号稳定：

\[
E[(a^l)^2]=E[(a^{l-1})^2]
\]

\Needspace{5\baselineskip}
可得：

\[
E[(a^{l-1})^2]
=\frac{1}{2}\cdot n_{\mathrm{in}}\cdot \mathrm{Var}(W^l)\cdot E[(a^{l-1})^2]
\]

得到W的方差应为2/n\_in.

\Needspace{5\baselineskip}
反向传播：

\Needspace{5\baselineskip}
设sum是L对第l层输入a的梯度（包含了第l层n\_out个神经元传来的梯度之和），f'(z)是激活函数的梯度，则由独立随机变量乘积方差公式：

\[
\mathrm{Var}\left(\frac{\partial L}{\partial z_j^{l-1}}\right)
=\mathrm{Var}(\mathrm{sum})\mathrm{Var}(f')
+\mathrm{Var}(\mathrm{sum})(E[f'])^2
+\mathrm{Var}(f')(E[\mathrm{sum}])^2
\]

\Needspace{5\baselineskip}
假设 \(E[\mathrm{sum}]=0\)：

\[
\mathrm{Var}\left(\frac{\partial L}{\partial z_j^{l-1}}\right)
=\mathrm{Var}(\mathrm{sum})\left(\mathrm{Var}(f')+(E[f'])^2\right)
\]

\Needspace{5\baselineskip}
即：

\[
\mathrm{Var}\left(\frac{\partial L}{\partial z_j^{l-1}}\right)
=\mathrm{Var}(\mathrm{sum})\cdot E[(f')^2]
\]

\Needspace{5\baselineskip}
已经推出：

\[
\mathrm{Var}(\mathrm{sum})
=n_{\mathrm{out}}\cdot\mathrm{Var}(W)\cdot
\mathrm{Var}\left(\frac{\partial L}{\partial z}\right)
\]

\Needspace{5\baselineskip}
并且对 ReLU 有 \(E[(f')^2]=0.5\)，代入：

\[
\mathrm{Var}\left(\frac{\partial L}{\partial z^{l-1}}\right)
=\left(n_{\mathrm{out}}\cdot\mathrm{Var}(W)\cdot
\mathrm{Var}\left(\frac{\partial L}{\partial z}\right)\right)\cdot\frac{1}{2}
\]

则W的方差应为2/n\_out.

一般来说，取前向传播对应的方差2/n\_in即可达到足够好的效果（一般用2/n\_in和2/n\_out初始化差别也不大）。

\FloatBarrier
\Needspace{5\baselineskip}
\hypertarget{ux53c2ux8003ux6587ux732e-9}{%
\subsection{参考文献}\label{ux53c2ux8003ux6587ux732e-9}}

\vspace{0.25em}
\begin{itemize}
\tightlist
\item
  Glorot, X., \& Bengio, Y. (2010). \href{https://proceedings.mlr.press/v9/glorot10a.html}{Understanding the difficulty of training deep feedforward neural networks}. AISTATS 2010.
\item
  He, K., Zhang, X., Ren, S., \& Sun, J. (2015). \href{https://arxiv.org/abs/1502.01852}{Delving Deep into Rectifiers: Surpassing Human-Level Performance on ImageNet Classification}. ICCV 2015.
\end{itemize}

\clearpage
